

\section{FAMILIAS DE FERMIONES}

{
\centering
\begin{tabular}{c|c|c}
\textbf{1ª familia} & \textbf{2ª familia} & \textbf{3ª familia} \\
\hline
$\binom{\nu_{eL}}{e_L^-}, e_R^-$ & $\binom{\nu_{\mu L}}{\mu_L^-}, \mu_R^-$ & $\binom{\nu_{\tau L}}{\tau_L^-}, \tau_R^-$ \\
$\binom{u_L}{d_L}, u_R, d_R$ & $\binom{c_L}{s_L}, c_R, s_R$ & $\binom{t_L}{b_L}, t_R, b_R$
\end{tabular}

\vspace{5pt}

\begin{tabular}{c|c|c|c|c|c|c|c}
\hline
& $u_R$ & $u_L$ & $d_R$ & $d_L$ & $\nu_{eL}$ & $e_R$ & $e_L$ \\
\hline
$Y$ & $4/3$ & $1/3$ & $-2/3$ & $1/3$ & $-1$ & $-2$ & $-1$ \\
\hline
$T_3$ & $0$ & $1/2$ & $0$ & $-1/2$ & $1/2$ & $0$ & $-1/2$ \\
\hline
$Q$ & $2/3$ & $2/3$ & $-1/3$ & $-1/3$ & $0$ & $-1$ & $-1$
\end{tabular}

\vspace{3pt}

}

Para la 2ª y 3ª familia es idéntico, con: $[u, \, d,\, e, \, \nu_e] \to [c, \, s, \, \mu, \nu_\mu] \to [t, \, b, \, \tau, \, \nu_\tau] $



\section{SIMETRÍA ELECTRODÉBIL (SM-EW): $SU(2)_L \times U(1)_Y$}
$\bullet$ \textbf{Quiralidad}: $\Psi_{L,R} = P_{L,R}\Psi$, con proyectores $P_{L,R} = \frac{1 \mp \gamma_5}{2}$.
\\ $\bullet$ \textbf{Generadores}: Para $SU(2)_L$, $T_i = \frac{\sigma_i}{2}$ (Isospín débil). Para $U(1)_Y$, hipercarga débil $Y$.
\\ $\bullet$ \textbf{Fermiones}: Agrupados en dobletes levógiros $L$ y singletes dextrógiros $R$.
\\ $\bullet$ \textbf{Carga eléctrica}: Relaciones: $Q = T_3 + \frac{Y}{2} \quad J_{\rm EM }^\mu = J_3 ^\mu + \frac{1}{2} J_{\rm Y}^\mu$.\\


$\bullet$ \textbf{Derivada covariante:} $D_\mu = \partial_\mu - igT_a W_\mu^a - ig'\frac{Y}{2}B_\mu$\\
$\bullet$ $W_\mu^a$ ($a=1,2,3$) bosones de gauge de $SU(2)_L$ con acoplo $g$. 
\\ $\bullet$ $B_\mu$ bosón de gauge de $U(1)_Y$ con acoplo $g'$.
\\ $\bullet$ \textbf{Bosones físicos}: Mezcla de $W_\mu^3$ y $B_\mu$ parametrizada por el ángulo débil $\theta_w$.
$$W_\mu^\pm = \frac{1}{\sqrt{2}}(W_\mu^1 \mp iW_\mu^2) \quad Z_\mu = W_\mu^3\cos\theta_w - B_\mu\sin\theta_w \quad A_\mu = W_\mu^3\sin\theta_w + B_\mu\cos\theta_w$$ 
$\bullet$ \textbf{Unificación}: $e = g \sin\theta_w = g' \cos\theta_w$. Para $Q^2= M_Z^2: e = 0.31, \, g = 0.65, \, g' =0.35$
\\ $\bullet$ \textbf{Masas}: Mecanismo de Higgs $M_Z = \frac{M_W}{\cos\theta_w}$; Constante de Fermi $\frac{G_F}{\sqrt{2}} = \frac{g^2}{8M_W^2}$. 

\section{LAGRANGIANO ELECTRODÉBIL}
\vspace{-2mm}
$$\mathcal{L}_{EW} = \sum_f \bar{\Psi}i\cancel{D}\Psi - \frac{1}{4}W_{\mu\nu}^a W_a^{\mu\nu} - \frac{1}{4}B_{\mu\nu}B^{\mu\nu}\;\;\;\;f \equiv \rm fermiones$$
$$W_{\mu\nu}^a = \partial_\mu W_\nu^a - \partial_\nu W_\mu^a + g\epsilon^{abc}W_\mu^b W_\nu^c \qquad B_{\mu\nu} = \partial_\mu B_\nu - \partial_\nu B_\mu$$ 

\section{REGLAS DE FEYNMAN (VÉRTICES)}
$\bullet$ \textbf{Corriente cargada ($W^\pm$)}: Interacciona solo con fermiones $L$. $W^\pm \rightarrow f \bar{f}': \;\; \frac{ig}{\sqrt{2}}\gamma^\mu P_L$
\\$\bullet$ Vértices con quarks: multiplica por elemento matriz CKM ($U_{ij}$) y tensor de color. 

$U_{\text{CKM}} = 
\begin{pmatrix}
|V_{ud} |& |V_{us}| & |V_{ub}| \\
|V_{cd}| & |V_{cs}| & |V_{cb} |\\
|V_{td}| & |V_{ts}| & |V_{tb}|
\end{pmatrix}
=
\begin{pmatrix}
0.974 & 0.224 & 0.00365 \\
0.224 & 0.973 & 0.0421 \\
0.00896 & 0.0413 & 0.999
\end{pmatrix}
\approx \mathbb{I} + 
\begin{array}{cc}
\text{correcciones}  \\
\text{pequeñas} 
\end{array}
$\\
$\bullet$ \textbf{Corriente neutra ($Z^0$)}: Interacciona con $L$ y $R$. $Z^0 \rightarrow f \bar{f}:\;\; \frac{ig}{\cos\theta_w}\gamma^\mu (g_L P_L + g_R P_R)$\\
$\bullet$ Vértices con neutrinos: multiplica por $\frac{1}{2}$ y sustituye la parte entre paréntesis por $P_L$.
\\$\bullet$ Vértices con quarks: multiplica por el tensor de color.

$\bullet$ \textbf{Acoplos quirales del $Z$}: $g_L = T_3 - Q\sin^2\theta_w$, $g_R = -Q\sin^2\theta_w$. 

\section{PROPAGADORES Y ESTADOS EXTERNOS}
$\bullet$ \textbf{Propagador fermión}: $\frac{i(\cancel{p}+m)}{p^2-m^2}$ (multiplicado por el tensor de color para quakrs). 
\\ $\bullet$ \textbf{Propagador bosón $V$ ($W, Z$)}: $\frac{-i(g_{\mu\nu} - k_\mu k_\nu/M_V^2)}{k^2-M_V^2}$. 
\\ $\bullet$ \textbf{Completitud}: Bosones masivos ($k^2 = M_V^2$) con 3 estados de polarización ($\lambda=\pm 1, 0$):
$$\sum_{\lambda} \epsilon_\mu^{(\lambda)}(k) \epsilon_\nu^{(\lambda)*}(k) = -g_{\mu\nu} + \frac{k_\mu k_\nu}{M_V^2} \epsilon_\mu(k) k^\mu = 0$$

\section{CÁLCULO DE AMPLITUDES (TRAZAS)}
$\bullet$ Amplitud cuadrada promediada: $|\overline{\mathcal{M}}|^2 = \frac{1}{\prod s_{in}} \sum |\mathcal{M}|^2$. 
\\ $\bullet$ \textbf{Propiedades trazas}: $P_L P_R = 0$, $P_L^2 = P_L$, $\gamma^\mu P_L = P_R \gamma^\mu$. 

$$\text{Tr}\{\gamma^\mu \gamma^\nu \gamma^\rho \gamma^\sigma\} = 4(g^{\mu\nu}g^{\rho\sigma} - g^{\mu\rho}g^{\nu\sigma} + g^{\mu\sigma}g^{\nu\rho}) \qquad \text{Tr}\{\gamma^\mu \gamma^\nu \gamma^\rho \gamma^\sigma \gamma_5\} = -4i\epsilon^{\mu\nu\rho\sigma}$$ 

$\bullet$ Contracción de tensores de Levi-Civita: $\epsilon_{\alpha\beta\mu\nu}\epsilon^{\rho\sigma\mu\nu} = -2(\delta_\alpha^\rho\delta_\beta^\sigma - \delta_\alpha^\sigma\delta_\beta^\rho)$ 

\section{CINEMÁTICA Y TASAS DE DESINTEGRACIÓN}
$\bullet$ Anchura parcial para desintegración $1 \rightarrow 2$: $\Gamma = \int \frac{1}{2M_W} |\overline{\mathcal{M}}|^2 d\text{Lips}(M_W; p_1, p_2)$\\
$\int d\text{Lips} = \frac{\sqrt{(M^2 - (m_1+m_2)^2)(M^2 - (m_1-m_2)^2)}}{8\pi M^2}$ \;\;si $m_1, m_2 \rightarrow 0 \Rightarrow \int d\text{Lips} = \frac{1}{8\pi}$.

$$d\text{Lips}_2 = (2\pi)^4 \delta^{(4)}(P - p_1 - p_2) \frac{d^3p_1}{(2\pi)^3 2E_1} \frac{d^3p_2}{(2\pi)^3 2E_2}\qquad d\text{Lips}_2 = \frac{|\mathbf{p}_f|}{16\pi^2 M} d\cos\theta \, d\phi$$

$\bullet$ \textbf{Sección eficaz}: $\frac{d\sigma}{d\Omega_{cm}} = \frac{1}{64\pi^2 s} |\overline{\mathcal{M}}|^2 \frac{|\mathbf{p}_f|}{|\mathbf{p}_i|}$. 


\section{ANCHURAS DE DESINTEGRACIÓN (W y Z)}
\textbf{Decaimiento del Bosón W:}
\\ $\bullet$ $\Gamma(W^- \to d\bar{u}) = 3 |V_{ud}|^2 \Gamma(W^- \to e^- \bar{\nu}_e) = 3 |V_{ud}|^2 ( g^2 M_W / 48\pi ) \;\;\; V_{ud} \overset{U_{CKM\approx 1}}{\approx} 1$
\\ $\bullet$ Factor de color: $F_c = \sum \delta_{ij}\delta_{ij}^* = 3$


\textbf{Decaimiento del Bosón Z:}
\\ $\bullet$ Anchura a neutrinos: $\Gamma(Z \to \nu_e \bar{\nu}_e) = g^2 M_Z/96\pi \cos^2\theta_w$ 
\\ $\bullet$ \textbf{Branching Ratios (BR)} aproximados: 
\\  Invisible ($Z \to \nu\bar{\nu}$): $\sim 20\%$ \;  Hadrónico ($Z \to q\bar{q}$): $\sim 70\%$ \; Leptónico ($Z \to l^+l^-$): $\sim 10\%$ 

\section{NARROW WIDTH APPROXIMATION (NWA)}
\textbf{Modificación del Propagador:} Para partículas inestables en canales $s$.
$$ \frac{1}{k^2 - M^2} \longrightarrow \frac{1}{k^2 - M^2 + i\Gamma M} $$

\textbf{Límite NWA ($\Gamma \ll M$):} Sustitución de la resonancia por la delta de Dirac.
$$ \frac{1}{(s - M^2)^2 + M^2 \Gamma^2} \longrightarrow \frac{\pi}{M\Gamma} \delta(s - M^2) $$ 
\\ $\bullet$ \textbf{Factorización de la Sección Eficaz:} Válido on-shell. $X$ partícula virtual. BR sig. pag. 

$$ \sigma(A+B \to C+D) \approx \sigma(A+B \to X) \cdot \text{BR}(X \to C+D) \quad \text{con } \sigma(A+B\to X) =\frac{\pi}{s}|\bar{F}|^2\delta (s-M_x^2)$$

\section{REGLAS DE FEYNMAN: BOSONES EW}
\textbf{Autointeracciones (Vértices de 3 y 4 puntos):}

% Generados por términos de Yang-Mills $-\frac{1}{4}W_{\mu\nu}W^{\mu\nu} - \frac{1}{4}B_{\mu\nu}B^{\mu\nu}$. % 
$\bullet$ Vértice $W^+ W^- \gamma$ $ -ie \left[ (k_1-k_2)_\lambda g_{\mu\nu} + (k_2-k_3)_\mu g_{\nu\lambda} + (k_3-k_1)_\nu g_{\lambda\mu} \right] $ 
\\ $\bullet$ Estructura tensorial vértices de 4 puntos: $ S_{\mu\nu,\lambda\rho} \equiv 2g_{\mu\nu}g_{\lambda\rho} - g_{\mu\lambda}g_{\nu\rho} - g_{\mu\rho}g_{\nu\lambda} $
\\ $\bullet$ Vértice $W^+ W^- \gamma \gamma$: $-i e^2 S_{\mu\nu,\lambda\rho}$ 

\section{BOSÓN DE HIGGS}
Escalar (Spin = 0), neutro, sin color. Propagador: $\frac{i}{p^2 - M_H^2}$ 

\textbf{Vértices de Interacción:} Los acoplamientos crecen con la masa de la partícula.
\\ $\bullet$ $H W^+ W^-$: $ig M_W g_{\mu\nu}$ \qquad $\bullet$ $H Z Z$: $\frac{ig M_Z}{\cos\theta_w} g_{\mu\nu}$ \qquad $\bullet$ $H f \bar{f}$ (Fermiones): $-i \frac{g}{2} \frac{m_f}{M_W} \delta_{ij}$ 

\textbf{Desintegración a Fermiones ($H \to f\bar{f}$):} 
\\ $\bullet$ Amplitud promediada (ej. $H \to b\bar{b}$): $|\overline{F}|^2 = 3 \frac{g^2 m_b^2}{4 M_W^2} \text{Tr}[(p_1+m_b)(p_2-m_b)]$ 
\\ $\bullet$ Anchura parcial dominante ($H \to b\bar{b}$): 
$$ \Gamma = \frac{3}{32\pi} \frac{g^2 m_b^2 M_H}{M_W^2} \left(1 - \frac{4m_b^2}{M_H^2}\right)^{3/2} $$ 
\\ $\bullet$ Branching Ratio principal: $\text{BR}(H \to b\bar{b}) \approx 60\%$ 

